\documentclass{ximera}

\input{../../preamble.tex}

\title{Inverse Trig Functions}

\begin{document}

\begin{abstract}
%
\end{abstract}
\maketitle




Sine, cosine, and tangent are functions. They connect angle measurements (degrees and radians) with ratios (numbers). Unfortunately, they are not one-to-one functions and do not have inverse functions. 

But, we would like them to.

In order to get inverse functions, we need to make sine, cosine, and tangent into one-to-one functions.  To accomplish this, we'll simply restrict their domains. 




We would like to analyze all of their characteristics and features. 


\begin{itemize}
     \item \textbf{\textcolor{red!70!black}{Domain}} 
     \item \textbf{\textcolor{red!70!black}{Zeros}} 
     \item \textbf{\textcolor{red!70!black}{Continuity}} \\
     \text{    }$\circ$ \textbf{\textcolor{purple!85!blue}{discontinuities}}\\
     \text{    }$\circ$ \textbf{\textcolor{purple!85!blue}{singularities}}\\
     \item \textbf{\textcolor{red!70!black}{End-Behavior}} 
     \item \textbf{\textcolor{red!70!black}{Behavior}} \\
     \text{    }$\circ$ \textbf{\textcolor{purple!85!blue}{intervals where increasing}}\\
     \text{    }$\circ$ \textbf{\textcolor{purple!85!blue}{intervals where decreasing}}\\
     \item \textbf{\textcolor{red!70!black}{Global Maximum and Minimum}} 
     \item \textbf{\textcolor{red!70!black}{Local Maximums and Minimums}} 
     \item \textbf{\textcolor{red!70!black}{Range}} 
     \item \textbf{\textcolor{blue!55!black}{...and we would like a nice graph}} 
\end{itemize}












\subsection*{Learning Outcomes}

\begin{sectionOutcomes}
In this section, students will 

\begin{itemize}
\item review inverse functions.
\item investigate inverse trigonometric functions.
\end{itemize}
\end{sectionOutcomes}









\begin{center}
\textbf{\textcolor{green!50!black}{ooooo-=-=-=-ooOoo-=-=-=-ooooo}} \\

more examples can be found by following this link\\ \link[More Examples of Inverse Trig Functions]{https://ximera.osu.edu/csccmathematics/precalculus/precalculus/inverseTrigFunctions/examples/exampleList}

\end{center}









\end{document}
